\documentclass[11pt]{article}
\usepackage[T1]{fontenc}
\usepackage{lmodern,amsmath,amssymb,amsthm,booktabs,array}
\usepackage[margin=1in]{geometry}
\usepackage[colorlinks=true,urlcolor=blue,citecolor=blue,linkcolor=blue]{hyperref}
\usepackage{xurl}
\newtheorem{theorem}{Theorem}
\newtheorem{lemma}[theorem]{Lemma}
\newtheorem{corollary}[theorem]{Corollary}
\newcommand{\scal}[1]{\langle #1\rangle_0}
\newcommand{\inv}{\iota}
\title{Permuted rotor loops beyond Jenzer (2026):\\
resolving Remark 5.5\\
(arXiv:2610.21593v1, p.~19)}
\author{Haobo Ma\\\small ChronoAI Pte Ltd, Singapore
\and Wenlin Zhang\\\small National University of Singapore\\\small ChronoAI Pte Ltd, Singapore}
\date{4 November 2026}
\begin{document}
\maketitle
\begin{abstract}
Extending Jenzer's (2026) same-order and star-family proofs, we resolve
his Remark 5.5 (arXiv:2610.21593v1, p.~19): every return permutation
satisfies the rotor-loop subset formula for every anticommuting unit
bivector family, without the star condition.
Conjugation eliminates every off-diagonal scalar contribution except
possibly complementary subsets when the family size is odd. Their total
vanishes because a balanced product in an exterior algebra has even degree.
Corollaries give mixed-sign compact and zero formulas, the exact chain
correction, and the range and separability of every return order.
For at most three factors the loop has only scalar and simple bivector
parts, and every scalar power is a Chebyshev polynomial of its scalar part.
\end{abstract}

\section{The question and the formula}
Jenzer's Remark 5.5 on page 19 of arXiv:2610.21593v1 suggests, on numerical
evidence, removing the star condition from Theorem 5.5~\cite{Jenzer}.
Theorem~\ref{main} proves that extension. This is the remark titled \emph{On the Scope, with
Numerical Evidence}; its source TeX label is \texttt{remark:5-7}.
The same-order case is already Jenzer's Theorem 5.2. The point here is
the cancellation for an arbitrary return order, including his triangle
example, where the star condition fails.
The earlier Zenodo version of 29 August 2026 labels the same posing
remark 5.7. Its identity-order formulas and two-factor power relation are
already proved there; the locator in the title refers to the arXiv PDF.
The star-family chain and same-order loop formulas appeared earlier in
Jenzer's \emph{Die Relationsalgebra}, released 6 August 2026,
Sections 3.1--3.2, S\"atze S5.10--S5.11, pp.~18--19~\cite{JenzerBand5}.

The cancellation ingredients have prior forms: conjugation averaging in
Shirokov~\cite[Theorem 3.2]{Shirokov}, and odd exterior cancellation in
Diehl--Reizenstein~\cite[Section 3.3, Lemma 22]{DR}, with a
determinant--Pfaffian antecedent in de Bruijn~\cite[Section 4]{deBruijn}.
Here they identify and cancel the complete complementary-subset contribution
for every return order.

Let $\mathcal C=\mathrm{Cl}(V,q)$ for a finite-dimensional real vector space
with nondegenerate quadratic form of any signature. Write
$t(x)=\scal{x}$ for its scalar part and $\widetilde x$ for reversion.
The blade and grade conventions agree with~\cite[Section 2.3]{LS}.
Suppose $k\geq1$ and the bivectors $B_1,\ldots,B_k$ satisfy
\begin{equation}\label{family}
B_i^2=-1,\qquad B_iB_j=-B_jB_i\quad(i\ne j).
\end{equation}
For $I=\{1,\ldots,k\}$ and $S=\{i_1<\cdots<i_m\}\subseteq I$, put
\[
B_S=B_{i_1}\cdots B_{i_m},\quad B_\varnothing=1,\qquad
c_i=\cos(\theta_i/2),\quad s_i=\sin(\theta_i/2),\quad w_i=s_i^2,
\]
\[
a_S=\prod_{i\in S}s_i\prod_{i\notin S}c_i,\quad
U_i=c_i+s_iB_i,\quad
\Theta_k(\tau)=(U_1\cdots U_k)
(\widetilde U_{\tau(1)}\cdots\widetilde U_{\tau(k)}).
\]
Here $\tau$ lists the return order, so that
\[
\inv(\tau,S)=\#\{(i,j)\in S^2:i<j,\ \tau^{-1}(i)>\tau^{-1}(j)\},
\qquad \eta(m)=(-1)^{m(m+3)/2}.
\]
The star condition means $t(B_S)=0$ for every nonempty $S$.

\begin{theorem}\label{main}
For every family satisfying \eqref{family}, every real parameter vector,
and every permutation $\tau$ of $I$,
\begin{equation}\label{formula}
\scal{\Theta_k(\tau)}=
\sum_{S\subseteq I}\eta(|S|)(-1)^{\inv(\tau,S)}
\prod_{i\in S}w_i\prod_{i\notin S}(1-w_i).
\end{equation}
No star condition is required.
\end{theorem}

\section{The cancellation argument}
The scalar map is tracial: $t(xy)=t(yx)$. Indeed, in an orthogonal basis,
the scalar part of a product of two distinct basis blades is zero, whereas
for equal blades the two products coincide. Bilinearity proves the assertion.
Consequently $t(gxg^{-1})=t(x)$ for invertible $g$.

\begin{lemma}\label{trace}
If $\varnothing\ne D\subseteq I$, then $t(B_D)=0$ unless $D=I$ and $k$
is odd. The assertion remains true if $B_i^2=\varepsilon_i\in\{-1,1\}$.
\end{lemma}
\begin{proof}
Conjugation by $B_j$ fixes $B_j$ and negates every other $B_i$, since
$B_j^{-1}=\varepsilon_jB_j$. Thus it multiplies $B_D$ by
$(-1)^{|D|-1}$ when $j\in D$, and by $(-1)^{|D|}$ when $j\notin D$.
For even nonzero $|D|$ choose $j\in D$; for odd $|D|$ with $D\ne I$
choose $j\notin D$. In either case invariance gives $t(B_D)=-t(B_D)$.
\end{proof}

\begin{lemma}\label{balanced}
In the exterior algebra on $e_1,\ldots,e_k$, the element
\[
W_\tau=(1+e_1)\cdots(1+e_k)
(1-e_{\tau(1)})\cdots(1-e_{\tau(k)})
\]
has no odd-degree component.
\end{lemma}
\begin{proof}
If $x,y$ are distinct exterior generators, then $x^2=y^2=0$ and $xy=-yx$.
Expansion of both sides gives
\begin{equation}\label{swap}
(1+x)(1+y)=(1+2xy)(1+y)(1+x).
\end{equation}
Every factor $1+2xy$ is central and has even degree. It is invertible,
with inverse $1-2xy$, since $(xy)^2=0$. Apply \eqref{swap} to adjacent
factors in the second word, with $x=-e_i$, $y=-e_j$, until their indices
are $k,k-1,\ldots,1$. This expresses the original second word as its
reverse-ordered version times a product $E$ of central even factors.
Now $(1+e_i)(1-e_i)=1$; cancelling from the middle out gives $W_\tau=E$.
Hence every homogeneous component of odd degree is zero.
\end{proof}

For comparison with~\cite[Lemma 22]{DR}, take the closed polygon with increments
\[
e_1,\ldots,e_k,-e_{\tau(1)},\ldots,-e_{\tau(k)}.
\]
Its signature maps to $W_\tau$ in the exterior algebra; the cited lemma,
applied to each odd coordinate projection, gives the same vanishing.
The odd-dimensional determinant formula in~\cite[pp.~138--139,
Eqs.~(4.3)--(4.5)]{deBruijn} supplies an older route: its extra
Pfaffian column is the sum of the increments, hence zero.

\begin{proof}[Proof of Theorem~\ref{main}]
Since reversion negates a bivector, $\widetilde U_i=c_i-s_iB_i$.
Expansion and sorting of the return factors give
\begin{equation}\label{expand}
t(\Theta_k(\tau))=
\sum_{S,T\subseteq I}a_Sa_T r_\tau(T)t(B_SB_T),
\qquad r_\tau(T)=(-1)^{|T|+\inv(\tau,T)}.
\end{equation}
Each swap of distinct generators contributes a minus sign; the number
of such swaps in the selected return factors is $\inv(\tau,T)$.

When $S=T$, moving equal generators together gives
$B_S^2=(-1)^{m(m+1)/2}$, where $m=|S|$: there are $m(m-1)/2$
interchanges and $m$ squares equal to $-1$.
The diagonal term in \eqref{expand} is therefore
$\eta(m)(-1)^{\inv(\tau,S)}a_S^2$, exactly the summand in \eqref{formula}.

For $S\ne T$, cancelling repeated indices expresses $B_SB_T$ as
$\pm B_{S\mathbin\triangle T}$. Lemma~\ref{trace} makes its scalar part
zero except possibly when $k$ is odd and $T=I\setminus S$.
For these remaining terms there are no repeated generators. Define
$\epsilon(S,T)$ by $e_Se_T=\epsilon(S,T)e_I$ when $T=I\setminus S$,
where all subset products are increasing. The same interchanges give
$B_SB_T=\epsilon(S,T)B_I$. Also
\[
a_Sa_{I\setminus S}=\prod_{i=1}^k s_ic_i.
\]
Thus the entire remaining off-diagonal sum is
\begin{equation}\label{correction}
\left(\prod_i s_ic_i\right)t(B_I)
\sum_{S\subseteq I}r_\tau(I\setminus S)\epsilon(S,I\setminus S).
\end{equation}
The last sum is precisely the coefficient of $e_I$ in $W_\tau$:
expanding its two words selects complementary subsets, and sorting the
second contributes $r_\tau$. For odd $k$ this coefficient is zero by
Lemma~\ref{balanced}. This proves \eqref{formula}. The proof divides by
no trigonometric coefficient, so zeros of $s_i$ or $c_i$ are included.
\end{proof}

\section{Consequences and source coverage}
\begin{corollary}\label{mixed}
Let $B_i$ be pairwise anticommuting bivectors with
$B_i^2=\varepsilon_i\in\{-1,1\}$. For arbitrary real $c_i,s_i$, put
$U_i=c_i+s_iB_i$. Then for every return order,
\begin{equation}\label{mixedformula}
t(\Theta_k(\tau))=
\sum_{S\subseteq I}(-1)^{|S|(|S|+1)/2+\inv(\tau,S)}
\prod_{i\in S}(\varepsilon_i s_i^2)\prod_{i\notin S}c_i^2.
\end{equation}
\end{corollary}
\begin{proof}
Lemma~\ref{trace} and the complementary-subset cancellation are unchanged.
The diagonal square becomes
$B_S^2=(-1)^{m(m-1)/2}\prod_{i\in S}\varepsilon_i$.
Multiplication by $(-1)^m$ from reversion yields \eqref{mixedformula}.
\end{proof}
With $c_i^2=1+\varepsilon_i s_i^2$, this is the arbitrary-permutation
extension of Jenzer's Theorem 6.10. The choice $\varepsilon_i=+1$ and
$(c_i,s_i)=(\cosh(\chi_i/2),\sinh(\chi_i/2))$ similarly extends his
Theorem 6.8. Both statements follow without the star condition.

\begin{corollary}\label{tracial}
Let $A$ be any unital associative real algebra with a linear functional
$t:A\to\mathbb R$ satisfying $t(1)=1$ and $t(xy)=t(yx)$.
If $B_i\in A$ anticommute pairwise and $B_i^2=\varepsilon_i\in\{-1,1\}$,
then \eqref{mixedformula} holds for
\[
(c_1+s_1B_1)\cdots(c_k+s_kB_k)
(c_{\tau(1)}-s_{\tau(1)}B_{\tau(1)})\cdots
(c_{\tau(k)}-s_{\tau(k)}B_{\tau(k)}).
\]
No dimension, grading or involution is needed for this algebraic statement.
\end{corollary}
\begin{proof}
The trace-support lemma uses only cyclicity and the inverses
$B_i^{-1}=\varepsilon_iB_i$. The diagonal reduction uses their scalar
squares and $t(1)=1$; the complementary-subset sum uses only the exterior
identity. These are exactly the steps proving Corollary~\ref{mixed}.
\end{proof}
This extends the observation in \cite[Remark 5.2, p.~17]{Jenzer} that
the same-order argument applies beyond bivectors. Taking $t$ to be
normalized matrix trace is one instance.

\begin{corollary}\label{mixedpolar}
For identity return order and normalized mixed-sign rotors
$c_i^2-\varepsilon_i s_i^2=1$, put $w_i=s_i^2$ and
\[
z_i=1+\varepsilon_iw_i+\mathrm i\varepsilon_iw_i
       =\rho_i e^{\mathrm i\alpha_i},\qquad
\rho_i=\sqrt{(1+\varepsilon_iw_i)^2+w_i^2},
\]
where $\alpha_i=\operatorname{atan2}(\varepsilon_iw_i,1+\varepsilon_iw_i)$.
Then
\begin{equation}\label{polar}
t(\Theta_k(\mathrm{id}))
=\operatorname{Re}\!\left[(1+\mathrm i)\prod_i z_i\right]
=\sqrt2\prod_i\rho_i\,
  \cos\!\left(\frac\pi4+\sum_i\alpha_i\right).
\end{equation}
It vanishes exactly when $\sum_i\alpha_i\in\pi/4+\pi\mathbb Z$.
Zeros exist for some parameters exactly when at least two signs
$\varepsilon_i$ agree. In particular a two-factor opposite-square loop has
scalar part $1+2w_1w_2$ and never vanishes.
\end{corollary}
\begin{proof}
The period-four sequence
$(-1)^{m(m+1)/2}=\operatorname{Re}[(1+\mathrm i)\mathrm i^m]$
turns \eqref{mixedformula}, with $\tau=\mathrm{id}$, into the product
in \eqref{polar}. Each $z_i$ is nonzero: for $\varepsilon_i=-1$,
$0\leq w_i\leq1$ and $\alpha_i\in[-\pi/2,0]$; for
$\varepsilon_i=+1$, $w_i\geq0$ and $\alpha_i\in[0,\pi/4)$.
Taking its real part proves the zero criterion, including elliptic
endpoints $w_i=1$. If two signs agree, set their weights to $1/\sqrt2$
and every other weight to zero; the two-factor scalar is
$1-2w_1w_2=0$. Without two equal signs there is at most one of each
type: one factor has scalar 1, while opposite types give $1+2w_1w_2$.
\end{proof}
For all negative signs, conjugating the product in \eqref{polar} recovers
Jenzer's compact form. The mixed-sign polar expression extends Sections
6.3--6.4 and the discussion after Theorem 6.10; using
$\operatorname{atan2}$ includes the endpoints omitted by the source's
half-open elliptic angle interval.

\begin{corollary}\label{chaincor}
Every mixed-sign family of Corollary~\ref{mixed} has the exact chain formula
\begin{equation}\label{chain}
t(U_1\cdots U_k)=
\begin{cases}
\prod_i c_i,&k\text{ even},\\
\prod_i c_i+t(B_I)\prod_i s_i,&k\text{ odd}.
\end{cases}
\end{equation}
The star condition is automatic unless $k\equiv3\pmod4$;
in that remaining case it is equivalent to $t(B_I)=0$.
\end{corollary}
\begin{proof}
Expand $\sum_S a_SB_S$ and apply Lemma~\ref{trace}. For odd $k$,
reversion gives $\widetilde{B_I}=(-1)^{k(k+1)/2}B_I$.
Since reversion fixes scalar parts, $t(B_I)=0$ also when
$k\equiv1\pmod4$. This describes the scalar obstruction, not a
classification of the underlying planes.
\end{proof}

For Jenzer's triangle in $\mathrm{Cl}_{0,3}$,
$B_1=e_1e_2$, $B_2=e_1e_3$, $B_3=e_2e_3$, one has $B_1B_2B_3=-1$.
Thus \eqref{chain} gives $c_1c_2c_3-s_1s_2s_3$, while
\eqref{formula} holds for all six permutations. For reverse order,
$\inv(\tau,S)=\binom{|S|}{2}$ makes every sign in \eqref{formula} positive,
so the sum is $\prod_i(w_i+1-w_i)=1$, agreeing with
$P_k\widetilde P_k=1$. Identity order recovers the already proved
Theorem 5.2 of~\cite{Jenzer}.
The triangle's scalar chain expression is already the classical
$XYZ$ quaternion formula in Henderson~\cite[Section 2.3, Eq.~(24), p.~10]{Henderson}.

\begin{corollary}\label{powers}
Let $1\leq k\leq3$, $B_i^2=\varepsilon_i\in\{-1,1\}$, and
$c_i^2-\varepsilon_i s_i^2=1$. For every return permutation,
$R=\Theta_k(\tau)$ has the form $R=a+K$, where $a=t(R)$ and $K$ is
a simple bivector or zero, with $K^2=a^2-1$. For every integer $n\geq0$,
\begin{equation}\label{cheb}
t(R^n)=T_n(a),\qquad T_0(x)=1,\quad T_1(x)=x,\quad
T_{n+1}(x)=2xT_n(x)-T_{n-1}(x).
\end{equation}
\end{corollary}
\begin{proof}
Reduce the expanded word using anticommutation and the scalar squares;
it is a linear combination of the $B_S$. For $k=3$ the coefficient of
$B_I$ is $\prod_i c_is_i$ times the top coefficient of $W_\tau$, hence
zero by Lemma~\ref{balanced}, independently of the signs of the squares.
All surviving subsets have size at most two. Products of two bivectors
have grades 0, 2 and 4; for $i\ne j$,
$\widetilde{B_iB_j}=B_jB_i=-B_iB_j$, so $B_iB_j$ is a bivector.
Thus $R=a+K$ with $K$ of grade two. Each factor has unit reversion norm,
so $1=R\widetilde R=a^2-K^2$. Its grade-four part gives
$K\wedge K=0$, the decomposability criterion for a bivector
\cite[Section 3.1, Example 3.1]{LS}. Finally $R^2-2aR+1=0$; multiplying
by powers of $R$ and taking scalar parts proves \eqref{cheb} by induction.
\end{proof}
This proves the three-factor observations in OF 2 and OF 3 of
\cite[Section 8.4, pp.~31--32]{Jenzer}, and extends the stated
two-factor power identity to all return permutations and mixed signs.
The Chebyshev consequence of the quadratic identity is classical; the
matrix-trace form appears in~\cite[Theorem 2, Eq.~(9)]{BR}.
The three-factor restriction matters. In $\mathrm{Cl}_{0,5}$, take
$B_i=e_0e_i$ for $1\leq i\leq4$ and $c_i=s_i=1/\sqrt2$.
Direct expansion of
$R=2^{-4}\prod_{i=1}^4(1+B_i)\prod_{i=1}^4(1-B_i)$ gives
$t(R)=-1/4$ and, writing $e_{ijkl}=e_ie_je_ke_l$,
\[
\langle R\rangle_4=\tfrac14(e_{0123}-e_{0124}+e_{0134}-e_{0234}+e_{1234}).
\]
For an even element in $\mathrm{Cl}_{0,5}$, the scalar of
$R\widetilde R=1$ is the sum of the squared blade coefficients.
Consequently
\[
t(R^2)=2t(R)^2-1+2\|\langle R\rangle_4\|^2=-1/4
\ne T_2(-1/4)=-7/8.
\]
The same scalar $-1/4$ occurs for a two-factor loop with
$(w_1,w_2)=(1,5/8)$, whose square has scalar part $-7/8$ by
Corollary~\ref{powers}; setting two further angles to zero realizes
that comparison within the same four-factor star family.

\begin{corollary}\label{graph}
For the square-minus-one formula, let $G_\tau$ be the graph on $I$ with
edge $\{i,j\}$ when $i<j$ and $\tau^{-1}(i)<\tau^{-1}(j)$.
Write $F_\tau(w)=t(\Theta_k(\tau))$ and let $e_\tau(S)$ count the
edges induced by $S$. Then
\begin{equation}\label{graphformula}
F_\tau(w)=\sum_{S\subseteq I}(-1)^{e_\tau(S)}
       \prod_{i\in S}w_i\prod_{i\notin S}(1-w_i).
\end{equation}
For reverse return order its range is $\{1\}$; for every other order
its range on $[0,1]^k$ is exactly $[-1,1]$, and it has a zero.
It factors across the connected components of $G_\tau$, and cannot
factor across any partition that cuts an edge. In particular it is a
product of single-variable functions exactly for reverse order.
\end{corollary}
\begin{proof}
The exponents $m(m+3)/2$ and $\binom m2$ have the same parity,
so the sign in \eqref{formula} is $(-1)^{e_\tau(S)}$.
The weights in \eqref{graphformula} are nonnegative and sum to one,
giving $|F_\tau|\leq1$. The graph is empty exactly for reverse order.
For an edge $\{i,j\}$, setting all other weights to zero gives
$1-2w_iw_j$: its endpoint values are $1,-1$, and
$(w_i,w_j)=(1,1/2)$ is a zero. Continuity supplies the whole interval.
Edge counts add over connected components, so their sums factor.
Conversely, any factorization across a partition cutting this edge
would restrict to $f(w_i)g(w_j)=1-2w_iw_j$ up to nonzero constants.
Setting either variable to zero forces both factors constant, a contradiction.
\end{proof}
This extends the same-order bound, zero existence and non-separability
statements in Sections 6.2--6.4 of~\cite{Jenzer} to every return order.
The tensor with vertex coefficients $(-1)^{e_G(S)}$ is also the
graph-state coefficient tensor of Hein--Eisert--Briegel~\cite[arXiv v1, Eqs.~(8)--(9)]{HEB}.
Their cut-rank formula~\cite[arXiv v1, Proposition 3, Eq.~(42)]{HEB} gives the
same criterion for factorization across a partition: no crossing edge.
The symmetric identity-order specializations in Section 6.6 are prior
results, not additional conjectures. In particular, if $a(n)$ denotes
OEIS A009545, then
\[
2^kF_{\mathrm{id}}(1/2,\ldots,1/2)
=\operatorname{Re}[(1-\mathrm i)(1+\mathrm i)^k]
=\operatorname{Im}(1+\mathrm i)^{k+1}=a(k+1).
\]
The entry gives $a(0)=0$, $a(1)=1$ and
$a(n)=2a(n-1)-2a(n-2)$~\cite{OEIS}. Thus this equal-weight sequence
is already tabulated independently of the rotor interpretation.

\section{Coverage and formal statements}
The source items below use rendered PDF numbering; several source TeX
labels differ. The scalar formula and the source's known special cases
are distinguished from the additional written corollaries.

\begin{center}\small
\begin{tabular}{@{}p{.30\linewidth}p{.64\linewidth}@{}}
\toprule Statement or source item & Coverage\\\midrule
Cancellation ingredients & Reproved; prior proofs Shirokov (2015), Diehl--Reizenstein (2018); determinant antecedent de Bruijn (1955).\\
Remark 5.5, p.~19 (TeX 5.7) & Proved: Theorem~\ref{main}, without the star condition.\\
Theorems 5.2, 5.5; reverse order & Reproved; prior proof Jenzer (2026): same order, the star case and reverse order.\\
Remark 5.2, p.~17 & Proved: arbitrary return order in a normalized tracial real algebra, Corollary~\ref{tracial}.\\
Theorem 4.2; Remark 4.2, p.~15 (TeX 4.4) & Proved: exact chain correction, Corollary~\ref{chaincor}; triangle expression reproved, prior formula Henderson (1977).\\
Theorems 6.8, 6.10, pp.~24--26 & Proved: arbitrary-permutation hyperbolic and mixed-sign extension, Corollary~\ref{mixed}.\\
Sections 6.3--6.4, 6.8 & Proved: mixed-sign compact and zero criterion, Corollary~\ref{mixedpolar}. Endpoint corrected ($w=1$): the half-open elliptic argument range.\\
Sections 6.2--6.4 & Proved: range and zeros for every return order, Corollary~\ref{graph}. Graph factorization reproved; prior proof Hein--Eisert--Briegel (2003; published 2004).\\
Section 6.6 & Reproved specialization; prior proof Jenzer (2026): symmetric and lacunary formulas. The half-weight sequence is OEIS A009545 shifted by one.\\
Section 8.4, OF 1 & Proved: scalar obstruction in Corollary~\ref{chaincor}. Not covered: classification of the planes.\\
Section 8.4, OF 2, pp.~31--32 & Proved for $k\leq3$: simple bivector part and no grade four, Corollary~\ref{powers}. Not covered: general higher-grade formulas.\\
Section 8.4, OF 3, p.~32 & Reproved; prior proof Jenzer (2026) for $k=2$ powers. Proved for $k=3$ powers, Corollary~\ref{powers}; four-factor obstruction exhibited. Not covered: general powers and other longer words.\\
\bottomrule
\end{tabular}
\end{center}

The linked Lean declarations cover the square-minus-one formula and its
scalar-cyclicity, trace-support and odd-coefficient ingredients.
At revision \nolinkurl{3bedf2870aacc6e7f17353444aadaa49fb73157c}, the declaration
\nolinkurl{D5.S3.QuadraticForms.JenzerPermutedRotorLoop.result} in
\href{https://github.com/the-omega-institute/trureturing/blob/3bedf2870aacc6e7f17353444aadaa49fb73157c/D5/S3/QuadraticForms/JenzerPermutedRotorLoop.lean}{\nolinkurl{JenzerPermutedRotorLoop.lean}}
proves the formula of Theorem~\ref{main}, with $\mathrm{Fin}\ k$ indexing
the family and the native Clifford scalar projection.
The companion
\href{https://github.com/the-omega-institute/trureturing/blob/3bedf2870aacc6e7f17353444aadaa49fb73157c/D5/S3/QuadraticForms/JenzerPermutedRotorLoopAlgebra.lean}{\nolinkurl{JenzerPermutedRotorLoopAlgebra.lean}}
contains \nolinkurl{NativeBridge.scalar_cyclic} and
\nolinkurl{TraceSupport.tracial_subset_vanishing}, expressing cyclicity
and Lemma~\ref{trace}'s square-$-1$ case. In
\href{https://github.com/the-omega-institute/trureturing/blob/3bedf2870aacc6e7f17353444aadaa49fb73157c/D5/S3/QuadraticForms/JenzerPermutedRotorLoopWords.lean}{\nolinkurl{JenzerPermutedRotorLoopWords.lean}},
\nolinkurl{ExteriorSum.signed_top_sum_zero} proves the odd top-coefficient
cancellation. The mixed-sign extension and \eqref{chain} are written
proofs above. The compact and zero formulas, chain refinement,
tracial-algebra extension, three-factor simplicity and powers, and graph factorization are also
written proofs, beyond the scope of the main declaration.

The accompanying \texttt{verify.py} compares exact polynomial coefficients
for every return permutation for $1\leq k\leq6$, checks all mixed
signatures through $k=4$, and checks the non-star triangle and a rational
rotation of it. These are finite checks with independent formal $c_i,s_i$;
the proofs above establish the all-parameter statements. It also checks
the mixed compact identity, graph coefficients and factorization,
three-factor grades and powers, and the four-factor example using exact
rational arithmetic; the ranges and counts are in \texttt{verification.txt}.

\begin{thebibliography}{99}\small
\bibitem{Jenzer} N.~Jenzer, \emph{Star-Ordered Rotor Loops in Clifford
Algebras}, arXiv:2610.21593v1 (2026), especially Theorem 5.5 and
Remark 5.5, pp.~18--19; Sections 6.7--6.8 and 8.4.
\url{https://arxiv.org/abs/2610.21593v1}. Earlier version:
29 August 2026, \url{https://doi.org/10.5281/zenodo.22159164}.
\bibitem{JenzerBand5} N.~Jenzer, \emph{Winkel statt Zeit -- Band 5 --
Die Relationsalgebra}, 6 August 2026, Sections 3.1--3.2,
S\"atze S5.10--S5.11, pp.~18--19.
\url{https://doi.org/10.5281/zenodo.20769239}.
\bibitem{Shirokov} D.~S.~Shirokov, \emph{Method of Averaging in Clifford
Algebras}, arXiv:1412.0246v2 (19 June 2015), Theorem 3.2, Eq.~(3.2), p.~4.
\url{https://arxiv.org/abs/1412.0246v2}.
\bibitem{DR} J.~Diehl and J.~Reizenstein, \emph{Invariants of multidimensional
time series based on their iterated-integral signature}, arXiv:1801.06104v2
(9 May 2018), Section 3.3, Lemma 22, p.~18.
\url{https://arxiv.org/abs/1801.06104v2}.
\bibitem{deBruijn} N.~G.~de Bruijn, \emph{On some multiple integrals involving
determinants}, Journal of the Indian Mathematical Society (N.S.) 19
(1955), 133--151, Section 4, pp.~138--139.
\url{https://pure.tue.nl/ws/files/1920642/597510.pdf}.
\bibitem{BR} P.~Brandi and P.~E.~Ricci, \emph{Composition Identities of
Chebyshev Polynomials, via $2\times2$ Matrix Powers}, Symmetry 12 (2020),
746, Theorem 2, Eq.~(9).
\url{https://doi.org/10.3390/sym12050746}.
\bibitem{HEB} M.~Hein, J.~Eisert and H.~J.~Briegel, \emph{Multiparty
entanglement in graph states}, arXiv:quant-ph/0307130v1 (18 July 2003),
Eqs.~(8)--(9), Proposition 3, Eq.~(42); published in Phys.~Rev.~A 69
(2004), 062311. \url{https://arxiv.org/abs/quant-ph/0307130v1}.
\bibitem{Henderson} D.~M.~Henderson, \emph{Euler Angles, Quaternions, and
Transformation Matrices---Working Relationships}, NASA-TM-74839,
JSC-12960 (July 1977), Section 2.3, Eq.~(24), p.~10.
\url{https://archive.org/download/NASA_NTRS_Archive_19770024290/NASA_NTRS_Archive_19770024290.pdf}.
\bibitem{LS} D.~Lundholm and L.~Svensson, \emph{Clifford algebra, geometric
algebra, and applications}, arXiv:0907.5356v1 (2009), Sections 2.3 and 3.1,
especially Example 3.1. \url{https://arxiv.org/abs/0907.5356v1}.
\bibitem{OEIS} OEIS Foundation Inc., \emph{The On-Line Encyclopedia of
Integer Sequences}, A009545, \emph{Expansion of e.g.f. $\sin(x)\exp(x)$},
formula for the imaginary part of $(1+\mathrm i)^n$ and its recurrence.
\url{https://oeis.org/A009545}. Accessed 4 November 2026.
\end{thebibliography}
\smallskip\noindent\small
\copyright\ 2026 Haobo Ma and Wenlin Zhang. Licensed under Creative
Commons Attribution 4.0 International (CC BY 4.0):
\url{https://creativecommons.org/licenses/by/4.0/}.
\end{document}
